\documentclass[12pt,a4paper]{article}
\usepackage[UTF8]{ctex}
\usepackage{geometry}
\usepackage{listings}
\usepackage{xcolor}
\usepackage{hyperref}
\usepackage{amsmath}
\usepackage{algorithm}
\usepackage{algpseudocode}

\geometry{margin=2.5cm}

\lstset{
    language=Python,
    basicstyle=\ttfamily\small,
    keywordstyle=\color{blue}\bfseries,
    commentstyle=\color{green!60!black},
    stringstyle=\color{red},
    numbers=left,
    numberstyle=\tiny\color{gray},
    stepnumber=1,
    numbersep=5pt,
    frame=single,
    breaklines=true,
    showstringspaces=false,
    tabsize=4
}

\title{Halo 系统完整代码技术文档}
\author{代码分析文档}
\date{\today}

\begin{document}

\maketitle

\tableofcontents
\newpage

\part{系统概述}

\section{文档概述}

本文档是 Halo 系统的完整技术文档，详细分析了系统的各个模块，从函数级别到语句级别解释了代码的作用和实现逻辑。

\section{系统架构}

Halo 系统采用分层架构：

\begin{enumerate}
    \item \textbf{配置层}：YAML 文件定义工作流
    \item \textbf{解析层}：Parser 将配置转换为 DAG
    \item \textbf{调度层}：Scheduler 生成执行计划
    \item \textbf{协调层}：Optimizer 管理多进程执行
    \item \textbf{执行层}：Worker 执行实际推理
\end{enumerate}

\section{系统工作流程}

\subsection{整体流程}

\begin{enumerate}
    \item \textbf{配置加载}：从 YAML 文件加载工作流定义
    \item \textbf{图构建}：Parser 解析配置，构建 Operator 对象图
    \item \textbf{调度规划}：根据选定的调度策略，生成执行计划
    \item \textbf{执行协调}：Optimizer 创建多个 Worker 进程，按计划派发任务
    \item \textbf{模型推理}：Worker 接收任务，执行批量推理
    \item \textbf{结果收集}：Optimizer 收集所有 Worker 的结果，更新 Query 状态
\end{enumerate}

\subsection{数据流}

\begin{verbatim}
YAML 配置
    ↓
Parser → Operator 图
    ↓
Scheduler → Workflows (执行计划)
    ↓
Optimizer → ExecuteInfo (任务封装)
    ↓
Worker → 推理结果
    ↓
Optimizer → 更新 Query.op\_output
    ↓
最终结果
\end{verbatim}

\part{Components 模块}

\section{Components 模块概述}

\texttt{components} 模块定义了 Halo 系统的核心数据结构，包括：
\begin{itemize}
    \item \texttt{Operator}：表示工作流中的一个操作节点
    \item \texttt{Query}：表示一个用户查询请求
    \item \texttt{ModelConfig}：模型配置信息
    \item \texttt{ExecuteInfo}：执行信息封装
\end{itemize}

\section{Operator 类}

\subsection{类定义位置}
\texttt{halo/components/operator.py}

\subsection{功能概述}
\texttt{Operator} 类表示智能体工作流图中的一个节点（OP），每个节点对应一个模型推理任务，包含模型配置、输入输出依赖关系等信息。

\subsection{关键属性}

\begin{lstlisting}
class Operator:
    def __init__(self, id=None, prompt=None, model_config=None, keep_cache=False):
        self.id = id if id is not None else uuid.uuid4()
        self.input_ops = []              # 前驱节点列表
        self.output_nodes = []           # 后继节点列表
        self.prompt = prompt             # 提示词（可选）
        self.model_config = model_config # 模型配置对象
        self.max_distance = None         # 到终点的最长距离
        self.keep_cache = keep_cache     # 是否保留 KV cache
\end{lstlisting}

\subsection{详细说明}

\subsubsection{\_\_init\_\_ 方法}
\begin{description}
    \item[功能] 初始化一个 Operator 对象
    \item[参数]
    \begin{itemize}
        \item \texttt{id}：操作符的唯一标识符，如果为 None 则自动生成 UUID
        \item \texttt{prompt}：可选的提示词字符串
        \item \texttt{model\_config}：\texttt{ModelConfig} 对象，包含模型相关配置
        \item \texttt{keep\_cache}：布尔值，指示是否保留 KV cache 供下游节点使用
    \end{itemize}
\end{description}

\subsubsection{运行时字段}
\begin{itemize}
    \item \texttt{data\_parallel}：标记该 OP 是否采用数据并行（多设备复制执行）
    \item \texttt{is\_duplicate}：标记是否为复制节点（用于数据并行）
    \item \texttt{duplicate\_info}：复制信息 \texttt{[rep\_idx, rep\_total]}
    \item \texttt{main\_node}：如果为复制节点，指向原始主节点
    \item \texttt{benchmark}：性能基准对象，记录 \texttt{init\_time}、\texttt{prefill\_time}、\texttt{generate\_time}
    \item \texttt{max\_distance}：由 parser 计算，表示到任意终点 OP 的最长路径长度
\end{itemize}

\subsection{Benchmark 类}

\texttt{Operator} 内部包含一个 \texttt{Benchmark} 类，用于记录性能指标：

\begin{lstlisting}
class Benchmark:
    def __init__(self):
        self.init_time = 0      # 初始化时间
        self.prefill_time = 0   # Prefill 阶段时间
        self.generate_time = 0  # Generate 阶段时间
    
    def total_time(self):
        return self.init_time + self.prefill_time + self.generate_time
    
    def update(self, dict):
        self.init_time += dict.get('init_time', 0.0)
        self.prefill_time += dict.get('prefill_time', 0.0)
        self.generate_time += dict.get('generate_time', 0.0)
\end{lstlisting}

\section{Query 类}

\subsection{类定义位置}
\texttt{halo/components/query.py}

\subsection{功能概述}
\texttt{Query} 类表示一个用户查询请求，包含查询 ID、提示词、执行状态、中间结果等信息。

\subsection{关键属性}

\begin{lstlisting}
class Query:
    def __init__(self, id, prompt, priority=0):
        self.id = id                    # 查询唯一标识
        self.prompt = prompt            # 用户输入的提示词
        self.prompt_len = len(prompt)   # 提示词长度
        self.status = "pending"         # 状态：pending/running/finished
        self.priority = priority        # 优先级（数值越大优先级越高）
        
        # 生成相关
        self.prefix_len = 0             # 前缀长度（用于 cache）
        self.input_ids = None           # Tokenized 输入 ID
        self.generated_tokens = 0        # 已生成 token 数
        self.decoded_tokens = ""        # 已解码的文本
        self.new_tokens = 0             # 本轮新生成的 token 数
        
        # 工作流执行相关
        self.op_output = {}              # 字典：op_id -> 该 OP 的输出文本
        self.step = 0                    # 当前执行步骤
        self.cache = None                # KV cache（用于 TransformersWorker）
        
        # 性能基准
        self.create_time = time.perf_counter()  # 创建时间戳
        self.benchmark = {}              # 字典：op_id -> (start_time, end_time)
\end{lstlisting}

\subsection{op\_output 字典}
这是 \texttt{Query} 类的核心字段之一，用于实现工作流中的依赖关系：
\begin{itemize}
    \item \textbf{键}：OP 的 ID（字符串）
    \item \textbf{值}：该 OP 对该查询产生的输出文本
    \item \textbf{用途}：当执行一个 OP 时，系统会检查其所有父 OP 的输出是否都在 \texttt{op\_output} 中，如果存在，则将这些输出拼接到当前 OP 的 prompt 后面，实现多步推理
\end{itemize}

\section{ModelConfig 类}

\subsection{类定义位置}
\texttt{halo/components/model\_config.py}

\subsection{功能概述}
\texttt{ModelConfig} 类封装了模型推理所需的所有配置参数。

\subsection{关键属性}

\begin{lstlisting}
class ModelConfig:
    def __init__(
        self, 
        model_name,              # 模型名称
        system_prompt=None,       # 系统提示词
        temperature=0.7,         # 采样温度
        top_p=0.9,               # Top-p 采样参数
        max_tokens=256,          # 最大生成 token 数
        max_batch_size=8,        # 最大批处理大小
        dtype='bfloat16',        # 数据类型
        use_chat_template=True,  # 是否使用聊天模板
        quantization=None,       # 量化配置
        lora_config=None,        # LoRA 配置
        max_model_len=None,       # 最大模型长度
        min_tokens=0,            # 最小生成 token 数
    ):
        # 所有参数直接赋值给实例属性
\end{lstlisting}

\section{ExecuteInfo 类}

\subsection{类定义位置}
\texttt{halo/components/exe\_info.py}

\subsection{功能概述}
\texttt{ExecuteInfo} 类是一个简单的数据容器，用于封装一次执行任务所需的所有信息。

\subsection{类定义}

\begin{lstlisting}
class ExecuteInfo:
    def __init__(self, op: Operator, query_ids, prompts):
        self.op = op                    # 要执行的 Operator
        self.query_ids = query_ids      # 查询 ID 列表
        self.prompts = prompts          # 对应的提示词列表
\end{lstlisting}

\part{Parser 模块}

\section{Parser 模块概述}

\texttt{parser.py} 模块负责从 YAML 配置文件解析并构建智能体工作流的有向无环图（DAG）。

\section{YAML 配置格式}

\subsection{基本结构}

\begin{lstlisting}
start_ops:
  - op0          # 起始节点列表
end_ops:
  - op2          # 终止节点列表

ops:
  op0:
    model: meta-llama/Llama-3.1-8B-Instruct
    prompt: "Please answer the question."
    max_tokens: 512
    input_ops: []           # 空列表表示起始节点
    output_ops:
      - op1
    
  op1:
    model: meta-llama/Llama-3.1-8B-Instruct
    prompt: "Please refine the answer."
    max_tokens: 1024
    input_ops:
      - op0      # 依赖 op0 的输出
    output_ops:
      - op2
\end{lstlisting}

\section{核心函数}

\subsection{build\_ops\_from\_config}

这是 parser 模块的核心函数，负责将配置字典转换为可执行的 \texttt{Operator} 对象图。

\subsubsection{执行流程}

\begin{enumerate}
    \item \textbf{验证必需字段}：检查 \texttt{ops}、\texttt{start\_ops}、\texttt{end\_ops} 是否存在且非空
    \item \textbf{创建 Operator 对象}：为每个 OP ID 创建一个空的 \texttt{Operator} 对象
    \item \textbf{验证引用和必需字段}：检查每个 OP 是否有 \texttt{model} 字段，引用的 OP ID 是否存在
    \item \textbf{建立依赖关系并配置模型}：
    \begin{itemize}
        \item 建立双向链接：\texttt{op.input\_ops} 和 \texttt{op.output\_ops}
        \item 自动推断 \texttt{keep\_cache}：如果后继节点使用相同模型，则自动设为 True
        \item 创建 \texttt{ModelConfig} 对象
    \end{itemize}
    \item \textbf{初始化运行时字段}：设置 \texttt{data\_parallel}、\texttt{is\_duplicate} 等
    \item \textbf{解析 start\_ops 和 end\_ops}：转换为 \texttt{Operator} 对象列表
    \item \textbf{计算最长距离}：调用 \texttt{\_compute\_max\_distances} 为每个 OP 计算到终点的最长路径
\end{enumerate}

\subsection{\_compute\_max\_distances}

\subsubsection{功能}

为每个 OP 计算到任意终点 OP 的最长路径长度（边数），同时检测循环依赖。

\subsubsection{算法：DFS + 记忆化}

使用深度优先搜索（DFS）和记忆化技术，时间复杂度 $O(V + E)$。

\begin{algorithm}
\caption{计算最长距离到终点}
\begin{algorithmic}[1]
\Procedure{DFS}{op, end\_set, memo, visiting}
    \If{op in memo}
        \State \Return memo[op]
    \EndIf
    \If{op in visiting}
        \State \textbf{raise} "Cycle detected"
    \EndIf
    \State visiting.add(op)
    \If{op in end\_set}
        \State memo[op] = 0
        \State visiting.remove(op)
        \State \Return 0
    \EndIf
    \State best = -1
    \For{child in op.output\_ops}
        \State d = DFS(child, end\_set, memo, visiting)
        \If{d != -1}
            \State best = max(best, d + 1)
        \EndIf
    \EndFor
    \State memo[op] = best
    \State visiting.remove(op)
    \State \Return best
\EndProcedure
\end{algorithmic}
\end{algorithm}

\part{Schedulers 模块}

\section{Schedulers 模块概述}

\texttt{schedulers} 模块实现了多种调度策略，用于将工作流图中的 \texttt{Operator} 分配到多个 GPU 设备上执行。

\section{调度器接口}

所有调度器都遵循相同的接口：

\begin{lstlisting}
def schedule_xxx(
    device_cnt: int,           # 设备数量
    start_ops: List[Operator],  # 起始节点列表
    end_ops: List[Operator],   # 终止节点列表
    all_ops: List[Operator],   # 所有 OP 列表
    queries: List[Query],      # 查询列表
    **kwargs                   # 其他参数
) -> List[List[Dict]]:
    """
    返回 workflows，格式为：
    List[ per-device List[ {"command": "execute", "params": (op, query_ids)} ] ]
    """
\end{lstlisting}

\section{调度策略详解}

\subsection{schedule\_rr：轮询调度}

最简单的调度策略：按照拓扑顺序，将每个 OP 轮询分配到设备上。

\begin{lstlisting}
def schedule_rr(
    device_cnt: int,
    all_ops: List[Operator],
    queries: List[Query],
) -> List[List[Dict]]:
    """Round-Robin per op; each op handles ALL queries."""
    workflows: List[List[Dict]] = [[] for _ in range(device_cnt)]
    all_ids = [q.id for q in queries]
    topo = _topo_order(all_ops)
    
    d = 0
    for op in topo:
        workflows[d].append({"command": "execute", "params": (op, list(all_ids))})
        d = (d + 1) % device_cnt
    return workflows
\end{lstlisting}

\subsection{schedule\_by\_levels：层级调度}

将 OP 按依赖层级分组，同一层级的 OP 可以并行执行。

\begin{lstlisting}
def schedule_by_levels(
    device_cnt: int,
    all_ops: List[Operator],
    queries: List[Query],
) -> List[List[Dict]]:
    """Level-based RR: assign ops within the same topo level round-robin."""
    workflows: List[List[Dict]] = [[] for _ in range(device_cnt)]
    all_ids = [q.id for q in queries]
    topo = _topo_order(all_ops)
    
    # 计算每个 OP 的层级
    level = {}
    for op in topo:
        level[op] = 0 if not op.input_ops else max(level[p] for p in op.input_ops) + 1
    
    # 按层级分配
    d = 0
    for lv in sorted(set(level.values())):
        layer_ops = [op for op in topo if level[op] == lv]
        for op in layer_ops:
            workflows[d % device_cnt].append({
                "command": "execute",
                "params": (op, list(all_ids))
            })
            d += 1
    return workflows
\end{lstlisting}

\subsection{schedule\_colocate\_parents：父节点共置调度}

优先将 OP 分配到其父节点所在的设备，以提高缓存局部性。

\subsection{schedule\_dp：数据并行调度}

将查询 ID 固定分片到各设备，每个 OP 在所有设备上并行执行。

\subsection{schedule\_heuristic：启发式调度}

结合多种优化技术的复杂调度策略：
\begin{itemize}
    \item 依赖感知的前沿扩展
    \item 基于 \texttt{max\_distance} 的优先级排序
    \item 数据并行复制（当设备数多于就绪 OP 数时）
    \item 父节点设备优先分配（缓存局部性）
\end{itemize}

\subsection{schedule\_search：Beam-Search 调度}

使用 Beam-Search 算法搜索近似最优的设备分配方案，考虑模型切换成本、负载均衡等因素。

\section{调度器对比}

\begin{table}[h]
\centering
\begin{tabular}{|l|c|c|c|c|}
\hline
调度器 & 并行执行 & 模型切换优化 & 缓存局部性 & 复杂度 \\
\hline
RR & 否 & 否 & 否 & 低 \\
\hline
By Levels & 是（层级内） & 否 & 否 & 低 \\
\hline
Colocate Parents & 否 & 否 & 是 & 中 \\
\hline
Data Parallel & 是 & 否 & 否 & 中 \\
\hline
Heuristic & 是 & 部分 & 是 & 高 \\
\hline
Beam-Search & 是 & 是 & 是 & 很高 \\
\hline
\end{tabular}
\caption{调度器特性对比}
\end{table}

\part{Optimizers 模块}

\section{Optimizers 模块概述}

\texttt{optimizers} 模块是 Halo 系统的核心协调器，负责：
\begin{itemize}
    \item 从 YAML 配置构建工作流图
    \item 创建和管理多个 Worker 进程（每个 GPU 一个）
    \item 调用调度器生成执行计划
    \item 协调 Worker 执行任务，保证依赖关系
    \item 收集结果并更新 Query 状态
\end{itemize}

\section{Optimizer 类}

\subsection{初始化}

\begin{lstlisting}
def __init__(self, config_path: str, **kwargs):
    self.config = load_config(config_path)
    mp.set_start_method("spawn", force=True)
    
    self.device_cnt = torch.cuda.device_count()
    self.processes: List[mp.Process] = []
    self.cmd_queues: List[mp.Queue] = []
    self.result_queues: List[mp.Queue] = []
    
    # 构建 OP 图
    self.ops, self.start_ops, self.end_ops, self.models = build_ops_from_config(self.config)
    self._create_workers()
\end{lstlisting}

\subsection{execute 方法}

这是 Optimizer 的核心方法，负责协调所有 Worker 执行任务。

\subsubsection{依赖关系保证机制}

\begin{itemize}
    \item \textbf{依赖检查}：在派发任务前，检查该任务的所有父 OP 是否已完成
    \item \textbf{阻塞处理}：如果依赖未满足，任务不会派发，等待父任务完成
    \item \textbf{自动重试}：每当有任务完成，会重新检查所有等待中的任务
    \item \textbf{死锁避免}：通过定期重试和依赖检查，避免系统死锁
\end{itemize}

\subsubsection{关键辅助函数}

\textbf{\_cmd\_transfer}：将调度器生成的 task 转换为 Worker 可执行的格式，拼接父节点的输出。

\textbf{\_task\_ready}：检查任务是否满足依赖关系，所有父 OP 的输出必须已就绪。

\textbf{\_try\_send}：尝试向 Worker 发送下一个任务，只有在依赖满足时才发送。

\part{Workers 模块}

\section{Workers 模块概述}

\texttt{workers} 模块实现了实际的模型推理执行层。Halo 支持两种 Worker 实现：
\begin{itemize}
    \item \texttt{vLLMWorker}：基于 vLLM 的高性能推理引擎
    \item \texttt{TransformersWorker}：基于 Transformers 库的推理引擎
\end{itemize}

\section{vLLMWorker 类}

\subsection{初始化}

\begin{lstlisting}
def __init__(
    self,
    id: int,
    physical_gpu_id: int,
    cmd_queue: "queue.Queue",
    result_queue: "queue.Queue",
) -> None:
    self.id = id
    self.device = f"cuda:{physical_gpu_id}"
    
    # 绑定到特定 GPU
    os.environ["CUDA_VISIBLE_DEVICES"] = str(physical_gpu_id)
    
    # 运行时状态
    self.model_name: Optional[str] = None
    self.llm: Optional[LLM] = None
    self.tokenizer: Optional[AutoTokenizer] = None
\end{lstlisting}

\subsection{init\_op 方法}

为每个 OP 初始化运行时状态，包括采样参数配置、模型切换（如果需要）、系统提示词设置。

\subsubsection{模型切换逻辑}

\begin{itemize}
    \item 比较当前模型名称和新 OP 的模型名称
    \item 如果不同，释放旧模型并加载新模型
    \item 清空 CUDA 缓存以释放显存
\end{itemize}

\subsection{execute 方法}

执行一个 OP 的批量推理任务，这是 Worker 的核心方法。

\subsubsection{执行流程}

\begin{enumerate}
    \item \textbf{初始化阶段}：调用 \texttt{init\_op} 初始化 OP 状态
    \item \textbf{输入构建}：根据是否使用聊天模板，构建输入
    \item \textbf{批量推理}：调用 \texttt{llm.generate(inputs, self.sampling\_params)}
    \item \textbf{结果处理}：提取每个查询的生成文本，记录性能基准
\end{enumerate}

\subsection{run 方法}

Worker 的主循环，持续监听命令队列并执行任务。

\part{Util 模块}

\section{Util 模块概述}

\texttt{util.py} 模块提供了一些通用的工具函数，主要用于：
\begin{itemize}
    \item 数据类型转换（字符串到 torch.dtype）
    \item GPU 设备管理（解析可见 GPU ID）
\end{itemize}

\section{函数详解}

\subsection{\_resolve\_dtype}

将字符串类型的数据类型规范转换为 \texttt{torch.dtype} 对象。

\subsection{\_visible\_physical\_gpu\_ids}

解析当前可见的物理 GPU ID 列表，考虑 \texttt{CUDA\_VISIBLE\_DEVICES} 环境变量的影响。

\part{使用指南}

\section{基本使用}

\begin{lstlisting}
from halo.optimizers import Optimizer
from halo.components import Query

# 创建优化器
opt = Optimizer("templates/adv_reason_3.yaml")

# 创建查询
queries = [
    Query(0, "What is Machine Learning?"),
    Query(1, "Explain neural networks."),
]

# 执行
opt.schedule(queries, strategy="heuristic")
queries = opt.execute(return_queries=True)

# 查看结果
for q in queries:
    print(f"Query {q.id} final output: {q.op_output}")

# 性能分析
opt.print_latency_percentiles()

# 清理
opt.exit()
\end{lstlisting}

\section{性能优化建议}

\subsection{调度策略选择}

\begin{itemize}
    \item \textbf{简单工作流}：使用 \texttt{rr} 或 \texttt{by\_levels}
    \item \textbf{多模型切换}：使用 \texttt{search} 或 \texttt{heuristic}
    \item \textbf{大批量查询}：使用 \texttt{dp}（数据并行）
    \item \textbf{缓存优化}：使用 \texttt{colocate\_parents}
\end{itemize}

\section{总结}

Halo 系统是一个完整的智能体工作流执行框架，具有以下特点：

\begin{itemize}
    \item \textbf{声明式配置}：通过 YAML 文件定义工作流，易于修改和扩展
    \item \textbf{多设备并行}：支持多 GPU 并行执行，提高吞吐量
    \item \textbf{智能调度}：多种调度策略，优化执行效率
    \item \textbf{依赖保证}：自动检查依赖关系，确保正确执行
    \item \textbf{高性能}：基于 vLLM，支持高效的批量推理
\end{itemize}

\end{document}
